\chapter*{勘误与增补}
\addcontentsline{toc}{chapter}{勘误与增补}
\SourcePage{185}
以下页码均指原书页面底部所印的页码。\footnote{原注 5。负行号表示从页底向上数。}
\begin{enumerate}[label=\arabic*.,leftmargin=2.5em,itemsep=.65\baselineskip]
\item 第 7 页，倒数第 15 行：将 \emph{normal (= à links connexes)}（正规，即链环连通）改为“余维 $\ge2$ 的单形处的链环”。
\item 第 8 页，第 6 行：\emph{en terme} 应为 \emph{en termes}。
\item 第 10 页，第 18 行：\term{perverse-sheaf}{反常层}的一个来源是 $\mathcal D$-模理论。由 [Kas75] 可知，复流形上的\term{holonomic-module}{完整模}的解\term{complex}{复形}与 de Rham \term{complex}{复形}在适当位移后都是\term{perverse-sheaf}{反常层}（按 [KS90] 定义 10.3.1 的意义），而且彼此\term{duality}{对偶}。
\item 第 19--20 页，1.1.4：

倒数第 5、4 行：将 \emph{foncteurs contravariants exacts à gauche}（左正合反变函子）改为 \emph{foncteurs additifs contravariants exacts à gauche}（左正合加性反变函子）。

关于\term{exact-category}{正合范畴}的讨论，有几处需要补充说明。尤其是第 20 页第 4--5 行的条件 (b) 表述不当，Gabber 给出了一个反例。

在这个反例中，令 $\mathcal A$ 为\term{torsion}{挠}子群被 $2$ 零化的有限生成阿贝尔群范畴（即形如 $\Z^n\oplus(\Z/2)^m$ 的群）。序列
\begin{equation*}\tag{$*$}
0\longrightarrow M\xrightarrow{i}N\xrightarrow{j}P\longrightarrow0
\end{equation*}
属于 $\mathcal E$，当且仅当它在阿贝尔群范畴中正合。此时，\term{admissible}{容许}满态射就是满同态；它们定义一个 Grothendieck 预拓扑，且 (b) 成立，但 $\mathcal A$ 并不正合，因为\term{admissible}{容许}单态射的复合 $\Z\xrightarrow{2}\Z\xrightarrow{2}\Z$ 不是\term{admissible}{容许}单态射。因此，Gabber 建议改为假设 $(\mathcal A,\mathcal E)^\sim$ 是\term{abelian-category}{阿贝尔范畴}，$A\mapsto h_A$ 有定义（即对正合序列 $(*)$，$j$ 是 $i$ 的余核），并补加条件：$A\mapsto h_A$ 的\term{essential-image}{本质像}对\term{extension}{扩张}稳定。注意，在 (a) 中，$F$ 应为加性函子；第 2 行 \emph{$F\mathcal A$ stable par extensions} 中应考虑的是\term{essential-image}{本质像}。在 (b) 中，$A\mapsto h_A$ \term{fully-faithful}{全忠实}只是 Yoneda 引理，但必须要求 $h_A$ 属于 $(\mathcal A,\mathcal E)^\sim$。

[要使 $(\mathcal A,\mathcal E)^\sim$ 成为\term{abelian-category}{阿贝尔范畴}，只需假设：
\SourcePage{186}
对 $\mathcal E$ 中的每个序列 $(*)$，$i$ 都是 $j$ 的核（于是 $F$ 属于 $(\mathcal A,\mathcal E)^\sim$ 的条件变为：对 $\mathcal E$ 中的每个序列 $(*)$，
$F(P)\xrightarrow{\sim}\ker(F(N)\rightrightarrows F(N\times_PN))$）；并且对每个\term{admissible}{容许}满态射 $N\to P$ 及态射 $Q\to P$，$N\times_PQ$ 可表，且 $N\times_PQ\to Q$ 同构于一个\term{admissible}{容许}满态射。]

注意，Quillen 所断言的 $\mathcal A$ 到 $(\mathcal A,\mathcal E)^\sim$ 的嵌入已在 [Lau83] 中证明。此外，还存在通常的大小问题：[在用集合与类奠基的范畴论中，当 $\mathcal A$ 不是小范畴时，无法构造函子 $\mathcal A^{\mathrm{op}}\to(\mathrm{Ab})$ 的范畴。]
\item 第 23 页，倒数第 1 行：TR 2 应为 TR 3。
\item 第 27 页，倒数第 2 行：“反之……”的断言需要 $0\in C$（更准确地说，$C$ 含有某个零对象）。反例：$\mathcal D=\D^b(k)$（$k$ 为域），$C$ 为非零向量空间范畴。
\item 第 31 页，顶部交换图下第 3 行：$\tau_{\le a}$ 应为 $\tau_{\ge a}$。
\item 第 31 页，倒数第 11 行：1.3.2 应为 1.3.2 (i)。
\item 第 32 页，倒数第 7 行：在 $(U,Z\ldots)$ 中，$Z$ 后应加逗号。
\item 第 33 页，第 14 行：\emph{a} 应为 \emph{à}。
\item 第 34 页，第 13--16 行：当 $I$ 为空区间（Bourbaki 允许这种区间）时，原定义使 $\D^{b,I}$ 成为空子范畴，但更自然的定义是令它为零对象所构成的范畴，从而对每个 $I$ 都有
\[
\D^{b,I}=\{A\in\D^b\mid H^jA=0\text{，对所有 }j\notin I\}.
\]
\item 第 35 页，倒数第 5、4 行：将 \emph{coreprésente}（\term{corepresent}{余表示}）改为 \emph{représente}（表示）。
\item 第 37 页，第 8 行行末：将 $T_1$ 改为 ${}^pT_1$。
\item 第 37 页，倒数第 13 行：$\tau_{\le0}T^*B$ 应为 $\tau_{\ge0}T^*B$。
\item 第 40 页，倒数第 3 行：删去“i.e.”前的字母 i。
\item 第 42 页，第 1、2 行：将 Walbroeck 改为 Waelbroeck [Wae72]。
\item 第 44 页，1.4.3 的第 3 行：将 $j_*$ 改为 $j^*$。
\item 第 44 页，(1.4.3.2) 的第 2 行：将 $j^*$ 改为 $j_*$。
\item 第 45 页，第 7 行：$j^*$ 应为 $j_*$。
\item 第 46 页，第 9 行：将 \emph{an} 改为 \emph{en}。
\item 第 53 页，(1.4.21.3) 上方第 2 行：删去 $T$。
\item 第 54 页，倒数第 4 行：$i_*H^0i^*B$ 应为 $i_*H^0i^*j_*B$。
\item 第 57 页，倒数第 1 行：$H^i{}^p\tau_{\le0}K$ 应为 $H^i{}^p\tau_{\ge0}K$。
\item 第 60 页，第 20、21 行：“有限维 $R$-向量空间的\term{local-system}{局部常值层}”这一说法略有不严；准确地说，是茎为有限维的\term{local-system}{局部常值层}。
\SourcePage{187}
\item 第 67 页，命题 2.2.2 (i)：应为 $\D_c^{\ge p}$（用小写 $p$，不用大写 $P$）。
\item 第 68 页，第 7 行：\emph{on} 应为 \emph{on a}。
\item 第 68 页，第 10 行：中间的 $A$ 应为 $i^*j_*A$。
\item 第 69 页，2.2.9：“在 $X$ 上，我们只考虑\term{etale}{平展}拓扑”这一说法不够严谨：\term{sheaf}{层}按\term{etale}{平展}拓扑考虑，而局部闭子集则是相对于 Zariski 拓扑而言的。另外，将 \emph{premier à}（与……互素）改为 \emph{différent de}（异于……）。
\item 第 70 页，倒数第 3 行：特别地，与\term{stratification}{分层}空间的情形（第 63 页）一样，有\term{self-dual}{自对偶}\term{perversity}{反常度}函数
\[
p_{1/2}:\mathcal S\longrightarrow\Z,\qquad S\longmapsto-\dim S.
\]
\item 第 71 页，第 2 行：这一表述引出了由\term{perversity}{反常度}函数构造 \term{t-structure}{$t$-结构}的推广，从而无需\term{constructible}{可构造}性，也无需使用\term{stratification}{分层}，见 [Gab04]。
\item 第 71 页，第 4 行：$<$ 应为 $>$。
\item 第 71 页，第 5 行：$>$ 应为 $<$。
\item 第 73 页，第 3 行：应为 $\D_c^b$（缺少上标 $b$）。
\item 第 73 页，第 7 行：应为 $\varprojlim{}^1\Hom^{-1}(K_n,L_n)$（缺少上标 $-1$）。
\item 第 79 页，3.1.5：最好假设 $\mathcal D$ 是严格全子范畴；否则，3.1.5 第 6 行的条件应为 $\gr_F^iK$ 属于 $\mathcal D$ 的饱和。
\item 第 80 页，倒数第 2 行：$(-i)^i$ 应为 $(-1)^i$。
\item 第 86 页，倒数第 9 行：注意，定理 3.2.4 蕴含各种背\term{site}{景}（凝聚情形等）中的\term{duality}{对偶}化对象都具有局部性。
\item 第 86 页，倒数第 5 行：\emph{unique}（唯一）应为 \emph{unique à isomorphisme unique près}（在唯一同构意义下唯一）。
\item 第 88 页，第 10 行：$\mathcal B_m$ 应为 $B_m$。
\item 第 97 页，命题 3.2.20 及其证明第 2、11 行：记号 ${}^p\D^{+\le0}$ 很可能应加逗号，写成 ${}^p\D^{+,\le0}$，与 \S 2.1 一致。（引理 3.2.21 及其证明也一样。）
\item 第 98 页，第 10 行：应为 ${}^p\D^{+,\le0}(X,\mathcal O)$、${}^p\D^{>0}(X,\mathcal O)$（加上左上标 $p$）。
\item 第 99 页，第 6 行：$R\Hom$ 应为 $R\underline{\Hom}$。
\item 第 99 页，倒数第 13 行：“\term{torsion}{挠}”的定义不够准确。应要求局部存在 $n\ne0$，使乘以 $n$ 的自同态为零。若再假设每个\term{stratum}{层片}都只有有限个连通分支，就无需再说“局部”。
\item 第 100 页，第 12 行：\term{l-adic}{$\ell$-进}形式体系的阐述见 [Eke90] 和 [BS15]；使用 $\infty$-范畴的较新阐述见 [GL17，第 II 章]。
\item 第 101 页，第 5 行：(2.2.12) 应为 (2.2.18)。
\item 第 101 页，第 8 行：2.2.12 应为 2.2.14。
\item 第 101 页，第 15 行：这里说只考虑系数 $\overline{\Q}_\ell$，但若干地方（4.1.13、4.3.1、4.3.2、4.3.3，以及 (4.2.5.1) 上方一行）使用的是系数 $\Q_\ell$。
\SourcePage{188}
\item 第 101 页，倒数第 8 行：$[p_{1/2}+]$ 应为 $[p_{1/2}^+]$。
\item 第 101 页，倒数第 6 行：$[p_{1/2}+]$ 应为 $[p_{1/2}^+]$。
\item 第 102 页，第 1 行：“参见 2.2.11”应为“参见 2.2.11 和 2.2.12”。
\item 第 103 页，倒数第 13 行：将 (4.1.7) 改为 (4.1.1)。
\item 第 114 页，倒数第 15 行：应为 ${}^p\D_c^{\le0}$，而非 ${}^p\D_c^{<0}$。
\item 第 119 页，第 4 行：$\widetilde{\operatorname{pr}}$ 应为 $\widetilde{\operatorname{pr}}_*$。
\item 第 120 页，推论 4.5.2 的证明：句末缺少句号。
\item 第 120 页，推论 4.5.4 第 2 行：应为 $\overline{\Q}_\ell$，而非 $\Q_\ell$。
\item 第 120 页，推论 4.5.5：应假设 $x$ 为闭点；否则，无法找到第 121 页头两行所说的 $f$。
\item 第 122 页，第 3 行：“我们只考虑分离\term{finite-type}{有限型}\term{scheme}{概形}……”——但在 5.4.8 中，还考虑了\term{strictly-henselian}{严格 Hensel} 化。
\item 第 123 页，公式 (5.1.2.2)：$R\Gamma.Ra_*$ 应为 $R\Gamma\,Ra_*$。
\item 第 123 页，倒数第 6 行：记号 $\widetilde{\Z}$ 通常写作 $\widehat{\Z}$（$\Z$ 的 profinite 完备化）。
\item 第 125 页，引理 5.1.3 (c)：缺少条件 $F\in P(Y_0)$。
\item 第 125 页，引理 5.1.3 和注 5.1.4：最好假设 $P(Y_0)$ 为严格全子范畴。
\item 第 128 页，5.1.13 第 4 行：将 $\D_{\le w}^b,\D_{\le w'}^b,$ 改为 $\D_{\le w}^b\times\D_{\le w'}^b$。
\item 第 128 页，倒数第 5 行：应为 \emph{d'après (i) et (ii)}（由 (i) 和 (ii)）。
\item 第 131 页，倒数第 5 行：$W^{w-1}$ 应为 $W_{w-1}$。
\item 第 131 页，倒数第 4 行：$U(X)$ 应为 $U\to X$。
\item 第 132 页，第 5 行：${}^p\varphi_*|F_0$ 应为 ${}^p\varphi_*F_0$。
\item 第 133 页，第一个交换图：$\mathfrak{p}r$ 应为 $\operatorname{pr}$。
\item 第 133 页，第一个交换图下第 3 行：$H_0\sim\mathbb A_0^{n-1}$ 应为 $H_0=\mathbb A_0^{n-1}$。
\item 第 134 页，第 3 行：“(1.3.17 (iii)”应为“(1.3.17 (iii))”。
\item 第 136 页，交换图下第 4 行：$V_0$ 应为 $v_0$。
\item 第 136 页，倒数第 10 行：“(5.1.8)”应为“(5.1.8))”。
\item 第 141 页，第 12 行：${}^P\tau_{>n}$ 应为 ${}^p\tau_{>n}$。
\item 第 142 页，倒数第 12 行：5.3.7 应为 5.3.8。
\item 第 143 页，第 11 行：在右括号“)”后插入逗号。
\item 第 143 页，第 16 行：$K=j_{!*}L[d][n]$ 应为 $K=(j_{!*}L[d])[n]$。
\item 第 146 页，倒数第 11 行：$\tau_{\le d-1}$ 应为 $\tau_{\ge d-1}$。
\item 第 148 页，第 5 行：在 6.1 的讨论中，除了从 $\F$ 过渡到 $\C$，还可以过渡到任意特征异于 $\ell$ 的代数闭域 $k$，并为 $k$ 上代数簇上具有\term{geometric-origin}{几何来源}的\term{semisimple}{半单} $\overline{\Q}_\ell$-\term{perverse-sheaf}{反常层}得到 6.2.5、6.2.8 和 6.2.10 的类似结论。
\SourcePage{189}
\item 第 149 页，第 11 行：Brian Conrad 在 2003 年注意到，由 $(A')$ 推出 $(A'')$ 的取极限论证需要使用\term{stratum}{层片}局部可缩这一事实。但 Gabber 指出，还需要在
\[
Z(\C)\xrightarrow{i}X(\C)\xleftarrow{j}U(\C)
\]
的情形下使用
\[
i^*j_*F=\varprojlim i^*j_*(F/\ell^nF),
\]
其中 $U$ 是若干\term{stratum}{层片}的开并，$Z$ 为其补集，$F$ 为 $U(\C)$ 上的\term{constructible}{可构造} $\Z_\ell$-\term{sheaf}{层}。对满足良好条件的\term{stratification}{分层}空间，例如来自局部有限三角剖分的\term{stratification}{分层}空间，这样做是可行的。
\item 第 153 页，第 14 行：\emph{nombre} 应为 \emph{nombres}。
\item 第 156 页，倒数第 14 行：这里还要求，对每个 $T\in T_S$，$L_S(T)$ 都是 $T$ 上\term{smooth}{光滑} $\Z/\ell\Z$-\term{sheaf}{层}的同构类组成的有限族，且这些\term{sheaf}{层}在几何纤维上不可约。
\item 第 157 页，第 7 行：$\xleftarrow{i^*}$ 应为 $\xrightarrow{i^*}$。
\item 第 157 页，第 8 行：使用离散赋值环并非必要；对任意夹在 $A$ 与 $\C$ 之间的\term{strictly-henselian}{严格 Hensel} 局部环 $V$，都有这样的\term{equivalence}{等价}。
\item 第 157 页，第 11 行：按假设，我们对于
\[
H^p\bigl(-,\underline{\Ext}^{q}(\mathcal F,\mathcal G)\bigr)
\]
具有比较同构，其中 $\underline{\Ext}$ \term{sheaf}{层}是相对于系数环 $\Z/\ell\Z$ 计算的。无法由此推出相对于 $\Z/\ell^n\Z$ 计算的 $\underline{\Ext}$ \term{sheaf}{层}的此类比较同构；不过，缩小 $S$ 后可以做到这一点，也可以利用标准\term{spectral-sequence}{谱序列}，推出相对于 $\Z/\ell^n\Z$ 计算的全局 $\Ext^p(\mathcal F,\mathcal G)$ 的比较同构。
\item 第 159 页，第 8 行：“在 $H^1(-,\Z/\ell)$ 上”应为“在 $H^1(-,\Z/\ell)$ 和 $H^2(-,\Z/\ell)$ 上”。
\item 第 160 页，倒数第 14 行：见 [Ham95]。
\item 第 160 页，倒数第 4 行：“可下降到 $S$”应为“可下降到某个适当的 $S$”。
\item 第 161 页，第 1 行：“因此在 $S$ 上仍成立”应为“因此在从 $S$ 中去掉 $\ell$ 为零的部分后仍成立”。
\item 第 161 页，第 7 行：这里的\term{spectral-sequence}{谱序列}与\term{truncation}{截断}单\term{pure}{纯}对象相联系，说明见 [MO15] 第 4.1.2 节。
\item 第 161 页，第 15 行：$Y$ 应为 $Y(\C)$。
\item 第 162 页，倒数第 4 行：“\term{finite-type}{有限型}\term{scheme}{概形}”应为“分离\term{finite-type}{有限型}\term{scheme}{概形}”。
\item 第 162 页，倒数第 3 行：将“若它属于满足以下条件的最小集合”改为“若 $(X,\mathcal F)$ 属于上述形状的对所构成的、对同构稳定且满足以下条件的最小类”。
\item 第 163 页，第 7 行：将“由此得到的各种稳定性，见”改为“特别地，对\term{intermediate-extension}{中间延拓}和\term{duality}{对偶}都稳定”。
\item 第 163 页，第 18 行：$\pi_*\C$ 应为 $\pi_*\underline{\C}$。
\item 第 163 页，倒数第 13 行：注意，\term{decomposition-theorem}{分解定理}在有限域上有更精确的版本：若 $f:X_0\to Y_0$ \term{proper}{固有}，$K\in\D_c^b(X_0,\overline{\Q}_\ell)$ 为\term{pure}{纯}反常对象，则 $Rf_*K$ 为\term{pure}{纯}反常对象的位移的有限直和；而 5.4.5 的分解在 $Y_0$ 上一般并不成立。利用
\SourcePage{190}
向射影情形的约化、\term{hard-lefschetz}{强 Lefschetz 定理} (5.4.10) 和\term{derived-category}{导出范畴}中的相关分解，这一事实早已为人所知，但直到最近才见于文献（也出现了更一般的形式），即 [dC15] 的命题 2.1.1 和 [SZ16] 的推论 3.2.5。
\item 第 164 页，第 10 行：$R\Hom$ 应为 $R\underline{\Hom}$。
\item 第 164 页，倒数第 11 行：$y\in Y$ 应为 $y\in V(\C)$。
\item 第 164 页，倒数第 7 行：$y\in Y$ 应为 $y\in Y(\C)$。
\item 第 165 页，注 6.2.11：关于这一问题的进展，例如可参见 M. Saito 的工作 [Sai88]、[Sai90]，T. Mochizuki 的工作 [Moc15]，以及 C. Sabbah 的 Bourbaki 报告 [Sab13]。
\item 第 167 页，第 14 行，[6]：第 77--79 页应为第 77--130 页。
\item 第 167 页，倒数第 4 行，[10]：\emph{Catégories abéliennes}（\term{abelian-category}{阿贝尔范畴}）应为 \emph{Catégories dérivées}（\term{derived-category}{导出范畴}）。
\end{enumerate}
